\section{Problem Definition}
An instance specifies shelf geometry, product dimensions and masses, supply
poses, target assignments, robot initialization, friction, and time limits.
The state contains robot position and velocity $(q,\dot q)$, product transforms
and velocities, contact information, and controller history. Initial experiments
use privileged simulator state. Perception uncertainty is a separate future
condition, not an implicit capability of this setup.

At control tick $t$, the policy supplies bounded native actuator references
$u_t$; physics advances at 2\,ms substeps with a nominal 20\,ms control period.
Products and the floating robot remain dynamic. Manipulation cannot modify
live poses, disable contacts, create hidden grasp attachments, or apply external
object forces. Reset fixtures and scratch-data predictions are separately
identified and excluded from physical success claims.

Let $S_i(x_t)$ mean that product $i$ lies inside its assigned shelf region,
is upright, has board contact, and has sufficiently small velocity. The current
environment requires these conditions continuously for a configured dwell time
$d$:
\begin{equation}
 G_{\rm shelf}(t) = \bigwedge_{i=1}^{N}
 \left[\forall\tau\in[t-d,t],\ S_i(x_\tau)\right].
\end{equation}
Stocking status is reversible: disturbing an earlier placement invalidates its
dwell. The default dwell is 0.5\,s. Assignment tolerances and complete geometry
checks are configuration-dependent and must accompany results.

The interaction protocol adds gates beyond the environment's geometric success:
verified opposing grasp contacts, actual lift without table support, stable
release, and sustained hand/product plus product/back contact for the final push.
Environment success alone does not establish these protocol events. Unassisted
success also requires zero logged operator interventions.

Recovery is a within-episode return from a specified failure to a valid task
state. We will record failure type, onset, attempts, recovery duration, and
whether recovery preserves earlier placements. Merely retrying, resetting, or
moving closer to a target does not constitute successful recovery.
